% SPDX-License-Identifier: CC-BY-4.0
% Copyright 2025 Florian Aram Feuerriegel — kassensturz.org
% Abschnitt 3 — Kalibrierung und Status quo (Modellwerte aus generiert/makros.tex)

\section{Kalibrierung und Status quo}
\label{sec:kalibrierung}

Die Bemessungsgrundlagen sind so gesetzt, dass das Modell im Jahr 2025 die
amtlichen Aufkommen trifft (Tabelle~\ref{tab:kalibrierung}). Steuern liegen
nahe an den Referenzwerten, die Beitragseinnahmen rund ein Achtel darüber;
Satzänderungen der Sozialversicherung wirken im Modell entsprechend etwas zu
stark. Die Ausgabenseite ist grob gegliedert. Ein Ausgleichsposten von \SQausgleich~Mrd.~\euro{}
für nicht modellierte Einnahmen, etwa Gebühren, Verkäufe, Vermögenseinkommen
und den Bundesbankgewinn, bringt den Saldo auf den Finanzierungssaldo der
Volkswirtschaftlichen Gesamtrechnungen (Tabelle~\ref{tab:ausgaben}). Die
Verteilungsmaße sind nicht kalibriert; sie folgen aus den Haushaltsdaten und
liegen nahe an den Werten von EU-SILC.

\begin{table}[t]
\centering\small
\caption{Kalibrierung im Jahr 2025. Aufkommen in Mrd.~\euro.}
\label{tab:kalibrierung}
\begin{tabular}{@{}lrrl@{}}
\toprule
 & Modell & Referenz & Quelle der Referenz \\
\midrule
Einkommensteuer einschl. Abgeltungsteuer & \KALest & 350 & BMF, Kassenjahr 2024 \\
Umsatzsteuer einschl. Einfuhrumsatzsteuer & \KALmwst & 290 & BMF, Kassenjahr 2024 \\
Körperschaftsteuer & \KALkst & 45 & BMF, Kassenjahr 2024 \\
Gewerbesteuer & \KALgewst & 75 & Destatis 2024 \\
Erbschaftsteuer & \KALerb & 12 & BMF, Kassenjahr 2024 \\
\COz-Bepreisung (BEHG, nationaler ETS-Anteil) & \KALco & 18 & DEHSt 2024 \\
Beiträge Rentenversicherung & \KALrv & 290 & DRV 2024 \\
Beiträge Krankenversicherung & \KALkv & 270 & BAS 2024 \\
\addlinespace
Finanzierungssaldo & \KALsaldo & $-$118,8 & Destatis, VGR 2024 \\
Gini-Koeffizient (Äquivalenzeinkommen) & \KALgini & 0,295 & EU-SILC 2024 \\
Armutsrisikoquote (\%) & \KALarmut & 15,5 & EU-SILC 2024 \\
\bottomrule
\end{tabular}
\end{table}

Ohne Politikänderung ist der Status quo nicht stabil
(Tabelle~\ref{tab:verlauf}). In der ersten Periode, 2025 bis 2028, liegt der
Saldo bei \SQsaldo~Mrd.~\euro{} je Jahr (\SQsaldoPct\,\% des BIP), weil die
Verteidigungsausgaben steigen und das Sondervermögen ausgibt. Der strukturelle
Saldo von \SQstruktur\,\% des BIP verfehlt die Grenze der Schuldenbremse von
$-0{,}70\,\%$. In der letzten Periode, 2041 bis 2044, treiben Alterung und
Zinsen das Defizit auf \SQsaldoPctEnde\,\% des BIP; die Schuldenquote steigt
von \SQschuldStart\,\%
auf \SQschuldEnde\,\%, und die kumulierten Emissionen überschreiten das
deutsche Budget für 1,7\,°C in den 2030er Jahren. Die Teams entscheiden also
nicht zwischen Beharren und Veränderung, sondern darüber, wer die Anpassung
trägt.

\begin{table}[t]
\centering\small
\caption{Status quo über fünf Perioden. BIP und Schuldenquote zu Beginn der
Periode, übrige Größen als Jahreswert der Periode; \COz{} kumuliert seit 2025
zu Periodenbeginn.}
\label{tab:verlauf}
% Erzeugt von docs/modell/messung.js — nicht von Hand bearbeiten
\begin{tabular}{@{}lrrrrrrrr@{}}
\toprule
 & BIP & Schulden & \multicolumn{3}{c}{Saldo} & Emiss. & CO\textsubscript{2} kum. & Gini \\
\cmidrule(lr){4-6}
Periode & Mrd.\,\euro{} & \% BIP & Mrd.\,\euro{} & \% BIP & strukt. & Mt/a & Mt & \\
\midrule
2025--2028 & 4.470 & 63,5 & $-$173,0 & $-$3,87 & $-$1,28 & 649 & 0 & 0,312 \\
2029--2032 & 4.968 & 71,5 & $-$240,7 & $-$4,85 & $-$1,43 & 500 & 2.373 & 0,312 \\
2033--2036 & 5.529 & 82,2 & $-$289,4 & $-$5,23 & $-$1,91 & 399 & 4.199 & 0,312 \\
2037--2040 & 6.143 & 93,4 & $-$291,5 & $-$4,75 & $-$2,25 & 314 & 5.669 & 0,312 \\
2041--2044 & 6.771 & 102,5 & $-$365,4 & $-$5,40 & $-$2,90 & 238 & 6.799 & 0,312 \\
\bottomrule
\end{tabular}

\end{table}
